% QMB_n09_algorithmen.tex — Chapter 9: Quantum Algorithms
% Substantive basis: Doc. 073, 147, 176, 316, 343; own presentation
\chapter{Quantum Algorithms: Deutsch, Grover, Shor and the Limits of the Period}
\label{ch:algorithmen}

\section{The Four Standard Algorithms}

Four algorithms form the textbook lesson of quantum computing. They show
what quantum interference can achieve and what it cannot achieve.

\subsection*{Deutsch: One Bit in One Query}

The Deutsch algorithm answers the question of whether an unknown
function $f:\{0,1\}\to\{0,1\}$ is constant or balanced, with a single
query --- classically one needs two.

In standard QM: One prepares $|0\rangle$, applies Hadamard, queries the oracle $U_f$, applies Hadamard, and measures. Result $0$:
constant; result $1$: balanced. The trick is phase kickback:
$U_f|x\rangle = (-1)^{f(x)}|x\rangle$. The Hadamard afterwards makes the
global phase difference visible.

In the FFGFT: The same sequence on the cylindrical phase space. The
oracle rotates the phase of the torsion vector; Hadamard swaps $z$ and
$r$ and rotates by $\pi/2$. The measurement result is identical. There is no
algorithmic difference --- only the conceptual framework is a
different one: The phase is not an abstract complex number, but the
rotational position of the torsion configuration.\status{K}

\subsection*{Grover: Search in $\sqrt N$ Steps}

Grover searches an unsorted database with $N$ entries in $O(\sqrt N)$
steps instead of $O(N)$ classically. The mechanism: One prepares the
register in a uniform superposition of all $N$ entries,
marks the sought entry by phase kickback
($U_\mathrm{Grover}:|x\rangle\to(-1)^{\delta_{x,x_0}}|x\rangle$) and
then reflects about the uniform superposition (inversion about
the mean). Each Grover step increases the amplitude of the sought
entry at the expense of all others. After $\pi/4\cdot\sqrt N$ steps,
the probability of the correct entry is near 1.

FFGFT classification: Grover iterations are, in the FFGFT, a sequence of
torsion rotations that systematically align the phase of the sought entry. The other entries land in unstable configurations;
the readout interaction drives the system deterministically into the
aligned pole.\status{K} The predictions are identical to
standard QM.

\subsection*{Shor: Factorization in Polynomial Time}

Shor's algorithm factors $N$ through period finding. Let $a$ with
$\gcd(a,N) = 1$; the function $f(x) = a^x\bmod N$ is periodic with
period $r$. The quantum Fourier transform (QFT) makes the period visible in
the measurement probabilities; continued fraction approximation extracts
$r$ from the measurement value. From $r$ the factors follow via $\gcd(a^{r/2}\pm1, N)$.

Central statement of the corpus: The algorithm decomposes into two processes
that have nothing to do with each other:\status{K}

\begin{center}
\small
\begin{tabular}{p{2cm}p{3cm}p{3cm}}
\toprule
 & \textbf{Preparation} & \textbf{Transformation} \\
\midrule
Task & Modular exponentiation $a^x\bmod N$ & Period extraction \\
Nature & arithmetic problem & wave operation (interference) \\
Cost & expensive ($O(n^3)$ gates) & cheap ($O(n^2)$ gates) \\
\bottomrule
\end{tabular}
\end{center}

The preparation accounts for $\approx95\,\%$ of the gate costs; the QFT only
$\approx5\,\%$.\status{K} Shor specified only the classical process
$(a,1)\to(a,a^x\bmod N)$ --- the implementation of modular
exponentiation in superposition is the actual engineering problem and
does not depend on quantum mechanical properties.

\section{$\phi$-QFT and Equivalence Theorem}

Doc.~147 develops an alternative Fourier transform that uses phases
$2\pi/\varphi^k$ instead of $2\pi/2^k$, where $\varphi$ is the
golden ratio:
\begin{equation}
  \phi\text{-QFT}:|x\rangle\mapsto\frac{1}{\sqrt{Q_\varphi}}\sum_{y=0}^{Q_\varphi-1}e^{2\pi ixy/Q_\varphi}|y\rangle,
\end{equation}
with $Q_\varphi = \varphi^n$ instead of $Q = 2^n$.

\textbf{Main theorem:} For Shor's algorithm for factorizing
$N < 2^{20}$ with error probability $\delta < 10^{-6}$:\status{K}
\begin{equation}
  P_{\mathrm{success}}(\text{Standard-QFT})\le P_{\mathrm{success}}(\phi\text{-QFT})\le P_{\mathrm{success}}(\text{Standard-QFT})+\xi.
\end{equation}
The $\phi$-QFT is equivalent to the standard QFT, with a maximum
deviation of $\xi\approx1{,}3\times10^{-4}$ --- negligible.
The equivalence theorem shows: $\xi$ has no structural function in $\mathbb{Z}/N\mathbb{Z}$. The period $r$ is determined by $\mathrm{lcm}(p-1,q-1)$,
an individual arithmetic property of $N$, not a
universal geometric parameter.

\section{The Structural Limits of the Period}

Three limits set precise bounds on what quantum algorithms can achieve.

\textbf{1. Weyl law and counting density.}\status{K} The eigenvalue density
of every compact state space grows according to a power law
$N(\lambda)\sim\lambda^{d/2}$. The prime counting function $\pi(T)$ and the
Riemann zeros grow like $T\ln T$. That is not a
power. A finite state space cannot fully carry such a spectrum. The Weyl
obstruction: An $n$-qubit register is a compact state space; arithmetic counting densities grow faster than any compact geometry allows.

\textbf{2. Preparation as bottleneck.}\status{K} The QFT makes
periodicities visible that are already present in the prepared states. The limiting step is state preparation
--- the evaluation of $a^x\bmod N$ --- not the transformation itself.
No geometric parameter $\xi$ can reduce the arithmetic complexity
of the preparation: It is a function of the prime structure of
$N$, not of the geometry of the carrier.

\textbf{3. No quantum advantage for RSA (as of 2026).}\status{X} An
advantage must become visible on the Shor factorization curve: The threshold of error correction must be demonstrated and maintained;
the preparation must be executed in superposition; the result
must be replicated with RSA-relevant key lengths (2048 bit).
None of this is realized as of 2026. Google Willow does surpass the error correction threshold for logical qubits --- but in
a regime far below the scaling necessary for Shor.

\section{Riemann Zeta and FFGFT Carrier}

Doc.~316 and Doc.~343 clarify the relationship between the Riemann zeta
and the FFGFT carrier.

\textbf{The Riemann zeta as a factor of the torus spectral zeta.}\status{K}
The closed form of the $\mathbb{Z}^4$ spectral zeta is:
\begin{equation}
  Z_{\mathbb{Z}^4}(s) = 8\,(1-4^{1-s})\,\zeta(s)\,\zeta(s-1).
\end{equation}
This is an identity. The zeros of the torus spectral zeta have
three sources: the Riemann zeros (from $\zeta(s)$), their shifted ones (from $\zeta(s-1)$) and the pole-zero $s = 1$. The
Riemann hypothesis is thus a statement about a geometric object
of the framework --- but one that the geometry cannot answer:
The torus can carry the zeros as a factor, but cannot generate them.

\textbf{$\mathbb{Z}_3$ orbifold projection.}\status{B} The
$\mathbb{Z}_3$ projection of $T^4/\mathbb{Z}_3$ transforms the full
Riemann zeta into a Dirichlet $L$-function:
\begin{equation}
  \zeta_{T^4/\mathbb{Z}_3}(s) = L(s,\chi_0) = (1-3^{-s})\,\zeta(s).
\end{equation}
The Euler product of this $L$-function is hierarchically ordered according to the entry depth
$k = \mathrm{ord}_p(3)$; the physically
visible primes $\{2,3,5,7,11,13\}$ are exactly those with $k\le6$.

\textbf{Galois structure and neutrinos.}\status{B} The four
Frobenius conjugacy classes of primitive elements in $\mathrm{GF}(27)^*$
are paired by the FFGFT sector pairing $k\leftrightarrow-k$ into two classes:
$\{f_3,f_4\}$ (Dirac, charged leptons) and $\{f_1,f_2\}$
(Majorana, $\nu_1,\nu_2$ active; $\nu_3$ in Dirac orbit $O_4$). Prediction:
$\nu_3$ has no Majorana phase.

\textbf{Weyl limit.}\status{K} No compact geometry can
carry the zeros of the Riemann zeta as a spectrum. The zero density
$N(T)\sim T\ln T$ demands a dilatation structure --- the unrolled
domain. The geometry shows the address, but does not provide the answer.
Three obvious geometric approaches are excluded: No
compact geometry provides the $T\ln T$ growth; $\xi$ does not shift the
critical line; the $\xi$ ladder does not discretize the zeros.

\section{Experimental Validation}

IBM Kingston (May 28, 2026), 50 independent runs, 2048 shots per run:
Bell state fidelity $0{,}9876\pm0{,}0028$. The repeatability variance
is statistically compatible with ordinary quantum shot noise (ratio
observed/expected variance: $1{,}02$). An earlier evaluation of 3
runs had erroneously reported $40\times$ determinism relative to QM; that was a sampling artifact.\status{K}

The $\xi$ correction to CHSH values is $\Delta\sim10^{-5}$ per
measurement, far below NISQ noise. The frequently cited cumulative
values for $N = 73$ qubits (Doc.~022 and Doc.~147) and the elementary
deviation $\xi/(2\pi)$ are not directly comparable and must not be
equated (Chapter~\ref{ch:chsh_exp}).

\quelle{Doc.~073 (Algorithms in T0, 2025), Doc.~147 §3 ($\phi$-QFT equivalence), §4 (Bell tests), §5 (Shor), §8 (IBM validation, May 2026), Doc.~176 (Preparation vs.\ transformation), Doc.~316 (Riemann, Weyl), Doc.~343 (Zeta Galois)}